% Epilogue of Deep Learning as an Experimental Science.
\specialchapter{Epilogue}{The laboratory notebook}{16}\label{ch:epilogue}
\begin{refsection}[chapters/16-epilogue.bib]
\labfloatprefix{E}

The fifteen chapters of this book examined other people's experiments. The epilogue runs one of its own, from the question to the report, and keeps its notebook in the open. The question is small on purpose: does label smoothing lower the test error of the small network that the laboratory sections have trained throughout? Label smoothing replaces the one-hot training target by a mixture of the one-hot vector and the uniform distribution over classes, with weight $0.1$ on the uniform part when it was introduced \parencite{SzegedyEtAl2016}, and its effects on accuracy, calibration, and distillation have been studied in image classification and machine translation \parencite{MullerKornblithHinton2019}. In \cref{ch:aggregates} it came out ahead of a baseline on every task of a suite, but there the two arms shared one fixed configuration and nobody asked how much of the gain survived a fair protocol. Every step below follows a protocol of an earlier chapter; the script that ran the whole study is \script{16-epilogue-study.py}, and it prints every number quoted here.

\addsec{The question and the predictions}

A question becomes an experiment when its unit, treatment, and response are fixed. The unit is a run of the network with two hidden layers of 64 units, trained for 30 epochs by stochastic gradient descent with momentum on the scikit-learn digits; the treatment is the smoothing parameter, $0$ or $0.1$; the response is the error on the 539 test images, measured once per run at the end of training. Before anything was run, four predictions were written into the header of the script (\cref{tab:epilogue:predictions}, middle column). They reached the repository in the same commit as the results, so nothing but this notebook vouches for their date; a commit before the first run would have. They are not forecasts to be proud of. They record what the experimenter expected, so that the experiment can show how far the expectation was from the result. The design below was fixed with them; four checks were added on review, after the first full run, and each is marked where it appears.

\begin{experiment}[the epilogue's study]
\field{question}{Does label smoothing ($0.1$) lower the test error of the digits network?}
\field{treatment}{smoothing parameter $0$ (baseline) or $0.1$}
\field{tuned}{learning rate and weight decay, on the same 10-point grid for both arms, 3 seeds per point, scored on a validation split carved from the training set}
\field{held fixed}{architecture, optimizer, epochs, batch size, data split, test set}
\field{varied}{initialization and data-order seeds, paired across the arms}
\field{response}{test error, read once per run}
\field{replicates}{chosen by a power calculation on a pilot}
\field{controls}{a positive control: the tuned baseline on half of the training set}
\field{analysis}{paired $t$-interval, sign-flip permutation test, and bootstrap for the difference; a $2\times2$ check of the interaction with weight decay}
\end{experiment}

\addsec{A pilot, and the number of seeds}

\Cref{ch:seeds} asks for the number of runs to be fixed before the comparison is made, and \cref{thm:toolkit:power} turns a pilot into that number. Ten paired seeds per arm at a learning rate of $0.05$, scored on a validation split of 300 images taken from the training set, gave validation errors of $3.97$ and $2.80$ percent and a standard deviation of the paired difference of $0.614$ points. To detect a difference of $0.25$ points with power $0.8$ at level $0.05$, the formula, with the variance of the paired difference in place of $2\sigma^2$, asks for $(z_{0.975}+z_{0.8})^2\,0.614^2/0.25^2\approx47.3$, so $48$ paired seeds. The pilot's own difference was large, but a pilot is for estimating the spread, not the effect; the effect it shows is a single noisy draw and was not used.

\addsec{Tuning with equal budgets}

\Cref{ch:budgets} showed how a comparison turns into a comparison of tuning budgets when the arms are not searched alike, and \cref{ch:ablations} how hyperparameters tuned for one arm favor it. Both arms were therefore searched over the same ten configurations, learning rates $0.01$, $0.02$, $0.05$, $0.1$, and $0.2$ with weight decay $0$ or $10^{-4}$, three seeds each, on the validation split only, ties going to the smaller learning rate and then the smaller weight decay. The baseline chose a learning rate of $0.05$ with weight decay $10^{-4}$, and label smoothing a learning rate of $0.1$ without weight decay, tied with two other configurations. At the baseline's choice, in a check added on review, label smoothing lowers the error by \pmval{0.86}{0.14} points (mean and $95$ percent $t$-interval over the paired seeds) instead of the $0.96$ reported below: the home-field advantage of \cref{ch:ablations} is about a tenth of the effect here.

A second check added on review concerns the grid. The baseline's weight decay sat on the edge of its axis, which \cref{prot:budgets:compare} rules out, so the axis was extended to $10^{-3}$ and $10^{-2}$ for both arms, on the validation split. The baseline's choice stayed ($3.67$ percent, against $3.78$ at $10^{-3}$, one validation image apart); label smoothing's best score was now shared by four configurations, and the tie rule chose a learning rate of $0.05$ with weight decay $10^{-3}$, where the effect is \pmval{-0.94}{0.12} points. The edge did not decide the comparison, but only the check could show that.

\addsec{Controls}

An experiment that cannot detect a known effect cannot be trusted to report the absence of an unknown one (\cref{ch:controls}). The positive control trains the tuned baseline on half of the training set. Its test error rose from $2.87$ to $4.27$ percent, a paired difference of \pmval{1.40}{0.16} points ($95$ percent $t$-interval over seeds): the pipeline sees an effect of this size easily. A negative control, a condition in which the effect cannot occur, was added on review: the tuned baseline rerun with $48$ fresh seeds came out \pmval{-0.12}{0.16} points from its confirmation runs (Welch's interval), an interval that contains zero, as it should.

\addsec{The result}

Forty-eight paired seeds per arm, each arm at its own tuned configuration, scored once on the test set, gave test errors of \pmval{2.87}{0.12} percent for the baseline and \pmval{1.91}{0.10} percent with label smoothing (means with $95$ percent $t$-intervals). The paired difference is $-0.96$ points with $95$ percent interval $[-1.10,-0.82]$; a bootstrap over seeds gives $[-1.09,-0.83]$, and a sign-flip permutation test of the paired differences gives $p<10^{-4}$ (none of 20{,}000 random sign flips produced a mean as far from zero). Label smoothing was better on $47$ of the $48$ seeds and tied on the remaining one (\cref{fig:epilogue:result}a).

These intervals are over seeds, for this split and these 539 test images: they say what new seeds would do, not what new digits would. \Cref{prot:seeds:measure} also asks for the noise of the test set, a check added here on review. All the errors of all the runs fall on 38 of the images; a baseline run misses $15.5$ of them on average and a run with label smoothing $10.3$, so the effect is about five images per run. A bootstrap that resamples the test images as well as the seeds \parencite{Owen2007} gives $[-1.65,-0.37]$. Label smoothing still wins, but by a margin known only to about $\pm0.6$ points, and more seeds would hardly narrow it; more test images would.

Pairing helped less than it often does: the paired interval is $0.14$ points wide on each side, against $0.16$ unpaired, because the outcomes of the two arms on the same seed correlate only weakly, $0.23$ ($95$ percent interval $-0.06$ to $0.48$). Their different learning rates are not the main reason. At a shared one, in the interaction check below, the correlation is $0.28$ without weight decay and $0.49$ with it ($95$ percent intervals $-0.01$ to $0.52$ and $0.24$ to $0.68$), against $0.95$ ($0.92$ to $0.97$) between baseline runs that differ only by a weight decay of $10^{-4}$: label smoothing changes which borderline digits a run gets wrong, and on this test set that is most of what a seed decides, the item churn of \cref{ch:seeds} ($82$ and $93$ percent of the seed variance of the two arms, with bootstrap intervals that start above $60$). Pairing costs nothing, and here it buys a little.

\begin{figure}[t]
\begin{fullwidth}
\centering
\begin{tikzpicture}
\begin{groupplot}[group style={group size=2 by 1,horizontal sep=20mm},lplot,width=0.45\linewidth,height=0.40\linewidth]
\nextgroupplot[title={(a) the paired seeds},xlabel={baseline test error (\%)},ylabel={label smoothing test error (\%)},xmin=1,xmax=4.6,ymin=1,ymax=4.6,minor tick num=1]
\addplot[reference,domain=1:4.6,samples=2] {x};
\addplot[run] table[x=baseline,y=smoothing] {figures/data/16-epilogue-paired.dat};
\nextgroupplot[title={(b) the $2\times2$ check},ylabel={test error (\%)},xlabel={weight decay},xmin=-0.5,xmax=1.5,xtick={0,1},xticklabels={$0$,$10^{-3}$},ymin=1.6,ymax=3.2,minor y tick num=1,legend pos=north east]
\addplot[armA,mark=*,error bars/y dir=both,error bars/y explicit,error bars/error bar style={inkblue,line width=0.5pt}] table[x=level,y=baseline,y error=baseline_hw] {figures/data/16-epilogue-cells.dat};
\addplot[armB,mark=square*,error bars/y dir=both,error bars/y explicit,error bars/error bar style={rubric,line width=0.5pt}] table[x=level,y=smoothing,y error=smoothing_hw] {figures/data/16-epilogue-cells.dat};
\legend{baseline,label smoothing}
\end{groupplot}
\end{tikzpicture}
\end{fullwidth}
\caption{The epilogue's study. (a) Each mark is one of $48$ seeds, trained in both arms at each arm's tuned configuration; points below the diagonal are seeds on which label smoothing had the lower test error. An error is a whole number of the 539 test images, so some marks coincide. (b) The $2\times2$ check at the baseline's tuned learning rate ($0.05$): mean test error over the same $48$ seeds with $95$ percent $t$-intervals, for weight decay $0$ or $10^{-3}$. Script \script{16-epilogue-study.py}.}\label{fig:epilogue:result}
\end{figure}

\addsec{An interaction check}

\Cref{ch:confounding,ch:ablations} warn that an effect measured at one setting of the other factors may not be the effect at another. A $2\times2$ design at the baseline's tuned learning rate crossed the smoothing parameter with weight decay $0$ or $10^{-3}$, again with $48$ paired seeds per cell (\cref{fig:epilogue:result}b). The main effect of label smoothing is \pmval{-0.80}{0.11} points ($95$ percent $t$-intervals over seeds), the main effect of weight decay \pmval{-0.16}{0.04}, and the interaction, the difference between the effect of smoothing with and without weight decay, \pmval{+0.15}{0.09}: smoothing helps a little less when weight decay is on ($-0.73$ points against $-0.88$), and helps in both cells. Two factors at two levels are not a characterization of the method, but they check the most common way for a single comparison to mislead, an effect that exists only at the settings where it was measured.

\addsec{Predictions against outcomes}

\Cref{tab:epilogue:predictions} sets the predictions beside what happened. Two of the four were right, the interaction only just, since its interval reaches $0.24$. The other two were wrong in the same direction. The effect was about twice the largest predicted, and the spread of the paired differences was larger than expected, so that the pilot asked for $48$ seeds where about $20$ had been guessed; nor did the predicted interval count the test images, which widen it to about $\pm0.6$ points however many seeds are run. A prediction that misses is not a failure of the experiment; it is the experiment's report on the experimenter. Had the predictions not been written down, the large effect would have been easy to remember as expected.

\begin{table}[t]
\centering\small
\renewcommand{\arraystretch}{1.15}
\begin{tabularx}{\linewidth}{@{}>{\raggedright\arraybackslash}p{0.27\linewidth}>{\raggedright\arraybackslash}X>{\raggedright\arraybackslash}X@{}}
\toprule
quantity & predicted before the runs & observed\\
\midrule
effect of label smoothing on test error & between $-0.5$ and $-0.1$ points & $-0.96$ $[-1.10,-0.82]$: larger\\
half-width of its 95\% interval & about $0.15$ points with about 20 seeds & $0.14$ with 48 seeds, $0.22$ projected for 20; about $0.64$ with the test images resampled\\
positive control (half the data) & worse by at least 1 point & worse by \pmval{1.40}{0.16}\\
interaction with weight decay & smaller than $0.2$ points & \pmval{+0.15}{0.09}\\
\bottomrule
\end{tabularx}
\caption{The predictions written into the study's script before the first run, and the outcomes, with $95$ percent intervals over seeds unless the test images are said to be resampled.}\label{tab:epilogue:predictions}
\end{table}

\addsec{The report}

A report states the claim with its uncertainty, the protocol that produced it, and what it does not show. For this study:

\begin{aside}[Report]
On the scikit-learn digits, label smoothing with parameter $0.1$ lowered the test error of a network with two hidden layers of 64 units from $2.87$ to $1.91$ percent: a paired difference of $-0.96$ points over $48$ paired seeds (better on $47$, tied on one), with $95$ percent interval $[-1.10,-0.82]$ on this test set and $[-1.65,-0.37]$ when its 539 images are resampled as well. Each arm was tuned over the same grid of ten configurations with three seeds each on a validation split, and widening the grid's weight decays moved the effect by $0.02$ points; the test set was read once per run. The number of seeds was fixed beforehand by a power calculation on a pilot. A positive control (half of the training set) raised the error by $1.40$ points, and a negative control (the baseline against itself) showed no difference. At the baseline's learning rate the effect of smoothing changed by \pmval{0.15}{0.09} points when weight decay was added. The resampled interval, the wider grid, and the negative control were added on review, after the first full run. The study does not show that the effect extends to other data, architectures, or training lengths, and it does not identify how smoothing helps; predictions written before the runs underestimated the effect by about half.
\end{aside}

The last sentence of the report is the one that most reports leave out. Its first half is \cref{ch:extrapolation} in one line; its second is \cref{ch:mechanisms}.

\addsec{Fifteen protocols}

Each chapter ended with a protocol. \Cref{tab:epilogue:protocols} lists them with the step that is most often skipped; the study above used six of them, those of \cref{ch:seeds,ch:budgets,ch:controls,ch:confounding,ch:ablations,ch:extrapolation}, and the predictions in its script are the first step of the last.

\begin{table}[t]
\centering\small
\renewcommand{\arraystretch}{1.12}
\begin{tabularx}{\linewidth}{@{}>{\raggedright\arraybackslash}p{0.36\linewidth}>{\raggedright\arraybackslash}X@{}}
\toprule
protocol & the step most often skipped\\
\midrule
\Cref{prot:seeds:measure}, \nameref{prot:seeds:measure} & Fix the number of runs before looking, and report every run.\\
\Cref{prot:instability:protocol}, \nameref{prot:instability:protocol} & Record every seed and version; a fixed seed buys an audit trail, not evidence.\\
\Cref{prot:aggregates:suite}, \nameref{prot:aggregates:suite} & Choose one aggregate and its interval in advance, and rerun every baseline under the same protocol.\\
\Cref{prot:budgets:compare}, \nameref{prot:budgets:compare} & Give the arms equal budgets over every hyperparameter that distinguishes them.\\
\Cref{prot:controls:search}, \nameref{prot:controls:search} & Run the placebo: random draws from the same space, through the same pipeline.\\
\Cref{prot:confounding:attribute}, \nameref{prot:confounding:attribute} & List every difference from the baseline, and cross the architecture with the recipe.\\
\Cref{prot:heldfixed:fixed}, \nameref{prot:heldfixed:fixed} & State the budget of the claim, and retune the hyperparameters that interact with the treatment at each of its levels.\\
\Cref{prot:ablations:design}, \nameref{prot:ablations:design} & Name the estimand, use a factorial design when components interact, and retune the ablated variant.\\
\Cref{prot:mechanisms:test}, \nameref{prot:mechanisms:test} & Define the mediator by a measurement, then knock it out and rescue it.\\
\Cref{prot:sharpness:sweep}, \nameref{prot:sharpness:sweep} & Give every configuration its own seeds, fix the measure first, and check it under rescaling.\\
\Cref{prot:fitting:fit}, \nameref{prot:fitting:fit} & Fix the objective in advance, start from a grid, and report convergence and bootstrap intervals.\\
\Cref{prot:protocol:card}, \nameref{prot:protocol:card} & Count every operation, and train each configuration as it would be trained at its budget.\\
\Cref{prot:extrapolation:before}, \nameref{prot:extrapolation:before} & Write the forecast and its interval down before the run finishes, and backtest the form.\\
\Cref{prot:emergence:announce}, \nameref{prot:emergence:announce} & Plot a continuous metric beside the headline one, with intervals and zero counts as bounds.\\
\Cref{prot:instruments:evaluate}, \nameref{prot:instruments:evaluate} & Sample the settings, score every arm under the same ones, and compare within settings.\\
\bottomrule
\end{tabularx}
\caption{The protocols of the fifteen chapters.}\label{tab:epilogue:protocols}
\end{table}

\addsec{What this book did not cover}

The book is about the training experiment, and three neighboring subjects were left to others. The science of benchmarks, of how a shared test set comes to rank models and what those rankings mean, is the subject of \textcite{Hardt2026}. Causal inference from observational data, where the experimenter cannot assign the treatment, is a field with its own textbooks \parencite{Pearl2009,ImbensRubin2015}; the book borrowed its potential outcomes for designed experiments (\cref{ch:confounding,ch:mechanisms}) and, for the correlations of a sweep (\cref{ch:sharpness}), only its first lesson, to compare within strata of the confounders. Human evaluation, where the instrument is a person, has methods of its own that \cref{ch:instruments} touched only through model judges. The documentation of data and models \parencite{GebruEtAl2021,MitchellEtAl2019} is part of reporting but was not treated here.

\addsec{The culture of experiments}

Every method in this book costs runs, and in a fast field the incentives reward a new result more than a checked one \parencite{SculleyEtAl2018}. Some of the remedies are institutional. Preregistration, the registration of a study's design and analysis before its outcomes are seen, is required by law for clinical trials in the United States, and services now offer it in every discipline \parencite{NosekEtAl2018}; predictions in the header of a script, committed before the first run, are its private version. Reproducibility programs and checklists at machine learning conferences ask authors to report the number of runs, error bars, and the computing infrastructure \parencite{PineauEtAl2021}. Other remedies are habits. Report negative results and failed replications as findings; a miniature that fails to reproduce part of a famous effect, as several in this book did, is information. Review a failed claim as an incident, without blame, as the prologue proposed. Keep the notebook.

The ornament at the foot of the epilogue's first page carries the sixteenth draw. The interval that ran off the page after two draws is now narrow, and the running mean, after a first draw $2.8$ standard deviations too high, sits $0.001$ below the line. Sixteen is still a small number of runs; but the estimate has an honest interval, which is more than one run could give, and that is the whole difference between a number and a measurement.

\notesandrefs

Label smoothing was introduced by \textcite{SzegedyEtAl2016} as a regularizer for the Inception networks, in which it lowered the top-1 and top-5 errors on ImageNet by about $0.2$ points. \Textcite{MullerKornblithHinton2019} found that besides helping accuracy it improves calibration, but that a network trained with it is a worse teacher for knowledge distillation, and traced the second effect to the penultimate layer, whose representations label smoothing draws into tight clusters, one per class. The study in this epilogue says nothing about those settings: it is a worked example of a protocol, not a contribution to the literature on label smoothing.

\textcite{NosekEtAl2018} argue the case for preregistration across the sciences and describe its forms, among them a data set split into a part for exploration and a sealed part for confirmation, which is what a validation split and a test set are. \textcite{PineauEtAl2021} report on the reproducibility program of NeurIPS 2019, including its checklist. Datasheets for datasets \parencite{GebruEtAl2021} and model cards \parencite{MitchellEtAl2019} are the standard proposals for documenting data and trained models. For crossed random effects \textcite{Owen2007} proposed the bootstrap that resamples rows and columns independently, here seeds and test images, and found it mildly conservative.

\printchapterbib
\end{refsection}
\labfloatrestore
